% ============================================================================
% Reinforcement Learning with Verifiable Rewards (RLVR)
% Section file: chapters/05_reinforcement_learning/03_rlvr.tex
%
% This file is designed to be included immediately after the corresponding
% \section command in the book's master file. It intentionally contains no
% \chapter, \section, package imports, or bibliography command.
% ============================================================================

Having established the actor--critic framework and the roles of the policy and
value functions, we can now examine how a language model can use an external
criterion of correctness to improve its generated responses. The central
question is how a scalar assessment of a completed response can be converted
into a learning signal for the token-level decisions that produced it. This
question motivates \emph{Reinforcement Learning with Verifiable Rewards}
(RLVR), a post-training approach in which the reward is obtained from a
verification procedure rather than exclusively from a learned model of human
preferences.

The distinction matters because not every desirable property of a language
model is equally difficult to evaluate. Helpfulness, style, and conversational
quality often require contextual judgments. By contrast, the final answer to a
well-specified mathematical problem can frequently be checked against a
known answer, and explicitly stated instruction constraints can sometimes be
tested programmatically. When a task admits a reliable verifier, the model can
be optimized against the outcome of its own generated response. The Tülu 3
training framework provides a concrete implementation of this idea for
mathematical reasoning and precise instruction following
\parencite{lambert2025tulu3}.

RLVR is best understood as a change in the source of the reward, not as a
replacement for the probabilistic generation mechanism of a language model.
The model remains an autoregressive policy, and its parameters are still
updated through reinforcement learning. What changes is the function that
evaluates a completed response. The following development begins with the
connection between autoregressive generation and episodic decision-making,
then formulates the verifiable reward objective, and finally explains how
PPO-based optimization, data construction, and the Tülu 3 experiments
instantiate the method.

% ----------------------------------------------------------------------------
% Main exposition
% ----------------------------------------------------------------------------

\subsection{From autoregressive language modeling to outcome-based learning}
\label{subsec:rlvr-autoregressive-foundation}

A language model represents a conditional probability distribution over
possible token sequences. Let \(x\) denote an input prompt and let
\(y=(y_1,y_2,\ldots,y_T)\) denote a response consisting of \(T\) generated
tokens. The model is parameterized by \(\theta\), and its autoregressive
factorization is

\begin{equation}
    \pi_\theta(y\mid x)
    =
    \prod_{t=1}^{T}
    \pi_\theta(y_t\mid x,y_{<t}),
    \qquad
    y_{<t}=(y_1,\ldots,y_{t-1}).
    \label{eq:rlvr-autoregressive-policy}
\end{equation}

The notation \(\pi_\theta\) emphasizes that the language model can be
interpreted as a stochastic policy. At each generation step, the model
assigns a probability to every token in its vocabulary, conditioned on the
prompt and the tokens already generated. A token is sampled or selected
according to this conditional distribution, appended to the sequence, and
used as part of the context for the next prediction. Generation ends when an
end-of-sequence (EOS) token is produced or when a prescribed generation limit
is reached.

This process has the structure of an episodic decision problem. In
reinforcement learning, an agent observes a state, selects an action according
to a policy, and moves to a subsequent state. For autoregressive generation,
the state at step \(t\) can be written as

\begin{equation}
    s_t=(x,y_{<t}),
    \qquad
    a_t=y_t.
    \label{eq:rlvr-state-action}
\end{equation}

The prompt and the previously generated tokens form the state, while the
next token is the action. Once the action is selected, the next state is
obtained by appending that token to the sequence. Conditional on the selected
token, this transition is deterministic at the level of the token history;
the stochasticity lies in the policy's selection of the token. A complete
response therefore corresponds to a trajectory through the state space of
the generation process.

The distinction between a token-level action and a sequence-level outcome is
central to RLVR. The verifier ordinarily evaluates the completed response,
not each token independently. A mathematical answer may be judged correct
only after the model has produced its final result, and an instruction may be
judged satisfied only after the complete response can be checked against its
constraints. Thus, a scalar reward can be assigned at the end of a trajectory
even though the policy made a sequence of token-level decisions.

In supervised fine-tuning, a model is trained to increase the likelihood of
a supplied reference response. If \(y^\star\) is the reference sequence, the
usual maximum-likelihood objective is to minimize the negative conditional
log-likelihood,

\begin{equation}
    \mathcal{L}_{\mathrm{SFT}}(\theta)
    =
    -\mathbb{E}_{(x,y^\star)}
    \left[
        \log \pi_\theta(y^\star\mid x)
    \right].
    \label{eq:rlvr-sft-objective}
\end{equation}

This objective directly rewards agreement with the demonstrated sequence.
It does not, by itself, distinguish between alternative responses that are
all correct. RLVR instead evaluates sampled responses according to whether
they satisfy a specified outcome criterion. Different reasoning trajectories
can consequently receive the same reward if they lead to an answer accepted
by the verifier.

The objective is therefore not to reproduce one prescribed token sequence.
It is to increase the probability that the policy generates a response
belonging to the set of responses accepted by the verification procedure.
This changes the source of supervision while preserving the autoregressive
structure of the language model.

\subsection{Verifiable rewards and the KL-regularized objective}
\label{subsec:rlvr-objective}

In conventional reinforcement learning from human feedback (RLHF), a learned
reward model commonly maps a prompt--response pair to a scalar score. Let

\begin{equation}
    r_\phi(x,y)\in\mathbb{R}
    \label{eq:rlvr-learned-reward}
\end{equation}

denote such a reward model, where \(\phi\) denotes its parameters. Its output
is an estimate of response quality learned from preference information.
Because it is a learned statistical model, its predictions can contain
estimation error and can reflect limitations in the preference data and the
reward-model training procedure.

RLVR replaces this learned estimate with a verification function. Let

\begin{equation}
    C(x,y)\in\{0,1\}
    \label{eq:rlvr-verifier}
\end{equation}

denote a binary verifier, where \(C(x,y)=1\) means that response \(y\)
satisfies the criterion associated with prompt \(x\), and \(C(x,y)=0\)
means that it does not. The corresponding reward is

\begin{equation}
    v(x,y)
    =
    \begin{cases}
        \alpha, & C(x,y)=1,\\
        0,      & C(x,y)=0,
    \end{cases}
    \label{eq:rlvr-binary-reward}
\end{equation}

where \(\alpha>0\) is the reward magnitude. Equivalently,

\begin{equation}
    v(x,y)=\alpha\,\mathbb{I}[C(x,y)=1],
    \label{eq:rlvr-indicator-reward}
\end{equation}

where \(\mathbb{I}[\cdot]\) is the indicator function. In the Tülu 3
implementation, the authors set \(\alpha=10\), based on pilot experiments,
and did not tune it further \parencite{lambert2025tulu3}.

A binary verifier maps a potentially complex response space to two reward
values. All responses that fail the specified condition receive zero,
regardless of how close they are to satisfying it. Responses that pass
receive the same positive reward, even if their reasoning paths or writing
styles differ. The reward therefore represents satisfaction of the
specified criterion, not a continuous estimate of general response quality.

For a fixed prompt \(x\), the expected verifiable reward under the current
policy is

\begin{align}
    \mathbb{E}_{y\sim\pi_\theta(\cdot\mid x)}
    [v(x,y)]
    &=
    \mathbb{E}_{y\sim\pi_\theta(\cdot\mid x)}
    \left[
        \alpha\,\mathbb{I}[C(x,y)=1]
    \right]
    \nonumber\\
    &=
    \alpha\,
    \Pr_{y\sim\pi_\theta(\cdot\mid x)}
    [C(x,y)=1].
    \label{eq:rlvr-expected-reward}
\end{align}

The equality follows from the definition of the expectation of an indicator
function. It gives the central interpretation of binary RLVR: maximizing
expected reward is equivalent to maximizing the probability that the policy
generates a response accepted by the verifier, up to the positive scaling
factor \(\alpha\).

Let \(\mathcal{D}\) denote the distribution of training prompts. Averaging
over this distribution gives

\begin{equation}
    J_{\mathrm{ver}}(\theta)
    =
    \alpha\,
    \mathbb{E}_{x\sim\mathcal{D}}
    \left[
        \Pr_{y\sim\pi_\theta(\cdot\mid x)}
        [C(x,y)=1]
    \right].
    \label{eq:rlvr-verification-objective}
\end{equation}

An unconstrained maximization of this objective can move the policy
substantially away from the distribution learned during pre-training and
earlier post-training. Such movement may improve the targeted verifier score
while degrading behavior that is not measured by that verifier. RLVR
therefore uses a regularization term based on the Kullback--Leibler (KL)
divergence between the current policy and a reference policy.

For discrete probability distributions \(P\) and \(Q\), the KL divergence is

\begin{equation}
    D_{\mathrm{KL}}(P\parallel Q)
    =
    \sum_z P(z)\log\frac{P(z)}{Q(z)}.
    \label{eq:rlvr-kl-definition}
\end{equation}

It is the expected logarithmic probability ratio under \(P\). The divergence
is nonnegative, provided the distributions satisfy the relevant support
conditions, and it is zero when the two distributions are identical. For a
prompt-conditioned language model, the relevant quantity is

\begin{equation}
    D_{\mathrm{KL}}
    \left(
        \pi_\theta(\cdot\mid x)
        \parallel
        \pi_{\mathrm{ref}}(\cdot\mid x)
    \right),
    \label{eq:rlvr-policy-kl}
\end{equation}

where \(\pi_{\mathrm{ref}}\) is the reference policy. The current policy
\(\pi_\theta\) is trainable, whereas the reference policy supplies the
distributional anchor against which the change in the current policy is
measured.

The KL-regularized RLVR objective used in Tülu 3 can be written as

\begin{equation}
    \boxed{
    \begin{aligned}
    \max_{\theta}\quad
    J_{\mathrm{RLVR}}(\theta)
    =
    \mathbb{E}_{x\sim\mathcal{D}}
    \Big[
        &\mathbb{E}_{y\sim\pi_\theta(\cdot\mid x)}
        [v(x,y)]\\
        &-\beta
        D_{\mathrm{KL}}
        \left(
            \pi_\theta(\cdot\mid x)
            \parallel
            \pi_{\mathrm{ref}}(\cdot\mid x)
        \right)
    \Big],
    \end{aligned}
    }
    \label{eq:rlvr-kl-objective}
\end{equation}

where \(\beta\geq 0\) controls the strength of the reference-policy
regularization. A larger \(\beta\) makes divergence from the reference policy
more costly; a smaller \(\beta\) permits greater movement in pursuit of the
verification reward. The coefficient therefore controls a trade-off between
targeted outcome optimization and preservation of the reference distribution.

The KL term can also be expressed at the token level. From the
autoregressive factorization in \cref{eq:rlvr-autoregressive-policy},

\begin{equation}
    \log \pi_\theta(y\mid x)
    =
    \sum_{t=1}^{T}
    \log \pi_\theta(y_t\mid x,y_{<t}).
    \label{eq:rlvr-log-sequence-prob}
\end{equation}

Applying the same identity to the reference policy gives

\begin{align}
    \log
    \frac{\pi_\theta(y\mid x)}
         {\pi_{\mathrm{ref}}(y\mid x)}
    &=
    \sum_{t=1}^{T}
    \log
    \frac{
        \pi_\theta(y_t\mid x,y_{<t})
    }{
        \pi_{\mathrm{ref}}(y_t\mid x,y_{<t})
    }.
    \label{eq:rlvr-token-kl-decomposition}
\end{align}

Taking the expectation under the current policy yields

\begin{equation}
    \begin{aligned}
    D_{\mathrm{KL}}
    \left(
        \pi_\theta(\cdot\mid x)
        \parallel
        \pi_{\mathrm{ref}}(\cdot\mid x)
    \right)
    =
    \mathbb{E}_{y\sim\pi_\theta}
    \left[
        \sum_{t=1}^{T}
        \log
        \frac{
            \pi_\theta(y_t\mid x,y_{<t})
        }{
            \pi_{\mathrm{ref}}(y_t\mid x,y_{<t})
        }
    \right].
    \end{aligned}
    \label{eq:rlvr-expected-token-kl}
\end{equation}

Thus, a sequence-level divergence can be understood as the expected sum of
token-level log-probability differences along trajectories sampled from the
current policy. The verifier evaluates the complete response, while the KL
regularizer accumulates differences between the conditional token
distributions of the current and reference policies.

The reward scale and the KL coefficient must be interpreted jointly. In
Tülu 3, a successful verification receives a reward of ten, while the KL
coefficient is varied across experiments. Changing the relative scale of
these terms changes the optimization trade-off even when the verifier itself
is unchanged.

The policy-gradient formulation explains why a nondifferentiable verifier
can still be used for learning. For a fixed prompt, define

\begin{equation}
    J(\theta)
    =
    \mathbb{E}_{y\sim\pi_\theta(\cdot\mid x)}
    [v(x,y)]
    =
    \sum_y \pi_\theta(y\mid x)v(x,y).
    \label{eq:rlvr-policy-gradient-start}
\end{equation}

Differentiating and applying the logarithmic derivative identity gives

\begin{align}
    \nabla_\theta J(\theta)
    &=
    \sum_y
    \nabla_\theta\pi_\theta(y\mid x)v(x,y)
    \nonumber\\
    &=
    \sum_y
    \pi_\theta(y\mid x)
    \nabla_\theta\log\pi_\theta(y\mid x)
    v(x,y)
    \nonumber\\
    &=
    \mathbb{E}_{y\sim\pi_\theta(\cdot\mid x)}
    \left[
        v(x,y)\nabla_\theta\log\pi_\theta(y\mid x)
    \right].
    \label{eq:rlvr-score-function-gradient}
\end{align}

The derivative is taken through the policy probability, not through the
verifier. The verifier may perform exact answer comparison or evaluate a
constraint using a discrete programmatic test; it does not need to be
differentiable with respect to \(\theta\). Using the autoregressive
factorization, the sequence-level gradient becomes

\begin{equation}
    \nabla_\theta J(\theta)
    =
    \mathbb{E}_{y\sim\pi_\theta}
    \left[
        v(x,y)
        \sum_{t=1}^{T}
        \nabla_\theta
        \log\pi_\theta(y_t\mid x,y_{<t})
    \right].
    \label{eq:rlvr-token-policy-gradient}
\end{equation}

A scalar outcome reward can therefore influence the parameters involved in
all token decisions along the sampled trajectory. For the full
KL-regularized objective, the gradient also contains the derivative of the
KL term. The score-function derivation above isolates the reward component
and shows why nondifferentiable outcome evaluation is compatible with
policy-gradient learning.

\subsection{PPO optimization, value estimation, and implementation}
\label{subsec:rlvr-ppo-implementation}

The policy-gradient estimator in \cref{eq:rlvr-score-function-gradient}
can have high variance because a single response-level reward is associated
with many token-generation decisions. Actor--critic methods address this
problem by learning a value function that estimates the expected return from
an intermediate state. This is the role of the value model in PPO-based
RLVR.

Let \(V_\psi(s_t)\) denote a value function parameterized by \(\psi\). It
estimates the expected return from state \(s_t=(x,y_{<t})\) under the policy
being considered. The value function is not the verifier and does not define
correctness. The verifier supplies the task reward; the value function
estimates expected returns and supports the construction of advantage
estimates.

The action-value function \(Q(s_t,a_t)\) represents the expected return
after taking action \(a_t\) in state \(s_t\). The advantage function is

\begin{equation}
    A(s_t,a_t)
    =
    Q(s_t,a_t)-V(s_t).
    \label{eq:rlvr-advantage-definition}
\end{equation}

An advantage is positive when the action performs better than the state-value
baseline and negative when it performs worse. In sampled training, an
estimated return \(\widehat{G}_t\) can be used to construct

\begin{equation}
    \widehat{A}_t
    =
    \widehat{G}_t-V_\psi(s_t).
    \label{eq:rlvr-advantage-estimate}
\end{equation}

Subtracting a state-dependent baseline can reduce the variance of the
policy-gradient estimator without changing its expectation, provided that
the baseline does not depend on the sampled action. In particular,

\begin{align}
    \mathbb{E}_{a_t\sim\pi_\theta(\cdot\mid s_t)}
    \left[
        \nabla_\theta\log\pi_\theta(a_t\mid s_t)
    \right]
    &=
    \sum_{a_t}
    \pi_\theta(a_t\mid s_t)
    \nabla_\theta\log\pi_\theta(a_t\mid s_t)
    \nonumber\\
    &=
    \sum_{a_t}
    \nabla_\theta\pi_\theta(a_t\mid s_t)
    \nonumber\\
    &=
    \nabla_\theta
    \sum_{a_t}\pi_\theta(a_t\mid s_t)
    =0.
    \label{eq:rlvr-baseline-zero}
\end{align}

It follows that subtracting a state-only baseline does not change the
expected policy gradient. A learned value function supplies such a baseline
while adapting it to the generation state.

A common estimator used with PPO is Generalized Advantage Estimation (GAE).
Define the temporal-difference residual as

\begin{equation}
    \delta_t
    =
    r_t+\gamma V_\psi(s_{t+1})-V_\psi(s_t),
    \label{eq:rlvr-td-residual}
\end{equation}

where \(r_t\) is the reward at step \(t\), \(\gamma\) is the discount factor,
and \(V_\psi\) is the value function. GAE combines residuals from multiple
future steps:

\begin{equation}
    \widehat{A}^{\mathrm{GAE}}_t
    =
    \sum_{l=0}^{T-t}
    (\gamma\lambda)^l\delta_{t+l},
    \label{eq:rlvr-gae}
\end{equation}

where \(\lambda\) controls the weighting of longer-horizon residuals. Lower
values of \(\lambda\) emphasize shorter-horizon estimates, while values
closer to one incorporate more of the later trajectory. The discount factor
\(\gamma\) controls the relative contribution of future rewards. Tülu 3 uses
\(\gamma=1.0\) and \(\lambda=0.95\) in its reported PPO configuration
\parencite{lambert2025tulu3}. With \(\gamma=1\), future rewards are not
geometrically discounted in this episodic formulation.

PPO updates the policy using trajectories sampled from a preceding policy.
Let \(\theta_{\mathrm{old}}\) denote the parameters of the policy that
generated the current rollout, and let \(\theta\) denote the parameters being
optimized. For a sampled state--action pair, define the probability ratio

\begin{equation}
    \rho_t(\theta)
    =
    \frac{
        \pi_\theta(a_t\mid s_t)
    }{
        \pi_{\theta_{\mathrm{old}}}(a_t\mid s_t)
    }.
    \label{eq:rlvr-ppo-ratio}
\end{equation}

A ratio greater than one means that the current policy assigns a higher
probability to the sampled action than the rollout policy did; a ratio below
one means that the probability has decreased. An unclipped surrogate
objective would use \(\rho_t(\theta)\widehat{A}_t\). Large changes in the
ratio, however, can produce overly aggressive updates and make the
optimization rely on samples generated by a substantially different policy.

PPO controls this incentive through the clipped surrogate objective,

\begin{equation}
    \boxed{
    L^{\mathrm{CLIP}}(\theta)
    =
    \mathbb{E}_t
    \left[
        \min
        \left(
            \rho_t(\theta)\widehat{A}_t,\,
            \operatorname{clip}
            \left(
                \rho_t(\theta),1-\epsilon,1+\epsilon
            \right)\widehat{A}_t
        \right)
    \right],
    }
    \label{eq:rlvr-ppo-clip}
\end{equation}

where \(\epsilon\) is the clipping coefficient. When the advantage is
positive, the objective favors increasing the probability of the sampled
action, but clipping limits the incentive for an excessively large increase.
When the advantage is negative, the objective favors decreasing that
probability, while clipping limits the incentive for an excessively large
decrease. Tülu 3 sets \(\epsilon=0.2\), giving the ratio interval
\([0.8,1.2]\) in the clipping operation.

The PPO clipping mechanism and the KL regularization in
\cref{eq:rlvr-kl-objective} should not be conflated. The ratio in
\cref{eq:rlvr-ppo-ratio} compares the current policy with the policy that
generated the present rollout. The KL term compares the current policy with
the reference policy. The two mechanisms use different reference
distributions and serve different purposes: clipping controls the local
policy update, whereas KL regularization anchors the policy to the selected
reference distribution.

The value model is trained to predict returns. A squared-error value loss
can be written as

\begin{equation}
    L_V(\psi)
    =
    \mathbb{E}_t
    \left[
        \left(
            V_\psi(s_t)-\widehat{G}_t
        \right)^2
    \right],
    \label{eq:rlvr-value-loss}
\end{equation}

where \(\widehat{G}_t\) is the return target used by the implementation.
The policy and value model therefore have distinct roles: the policy learns
which token actions should become more probable, while the value model
estimates expected returns to support advantage estimation.

The Tülu 3 implementation uses a value-function coefficient \(c_1=0.1\),
a gradient-norm threshold of \(1.0\), and a linear learning-rate schedule.
It also normalizes the advantage estimates by subtracting their batch mean
and dividing by their batch standard deviation:

\begin{equation}
    \widehat{A}^{\mathrm{norm}}_t
    =
    \frac{
        \widehat{A}_t-\mu_A
    }{
        \sigma_A
    },
    \label{eq:rlvr-advantage-whitening}
\end{equation}

where \(\mu_A\) and \(\sigma_A\) are the mean and standard deviation of the
advantages over the normalization batch. This operation controls the
numerical scale of the policy-gradient signal. The implementation sets
dropout to zero during reward-model and reinforcement-learning training so
that rollout-time and learning-time token log probabilities remain
consistent. If stochastic dropout causes these probabilities to differ,
the probability ratio can depart from one even at the start of the first PPO
epoch, leading to excessive clipping and potentially eliminating useful
gradients \parencite{lambert2025tulu3}.

The response termination condition is also treated explicitly. PPO samples
responses up to a maximum token length, and a response that reaches this
limit without an EOS token may be incomplete. Tülu 3 applies an additional
penalty of \(-10\) to responses without an EOS token. This termination
penalty is an implementation detail in addition to the binary verification
reward; it encourages the policy to complete its responses rather than
repeatedly generating to the maximum length.

The complete relationship between generation, verification, regularization,
value estimation, and policy updating is summarized in
\cref{fig:rlvr-training-pipeline}. The policy generates a response from a
prompt, the verifier evaluates the completed response, and the resulting
reward is used together with the KL term and the advantage estimates in the
PPO update. The reference policy supplies the distribution used for the KL
regularizer, while the value model supplies estimates used to construct
advantages. The verifier and value model are therefore not interchangeable:
one defines the task outcome, and the other supports optimization.

\begin{figure}[htbp]
    \centering
    % ============================================================================
% RLVR pipeline schematic.
% This file contains only the drawing code. The figure environment and its
% caption are defined in chapters/05_reinforcement_learning/03_rlvr.tex.
% Uses the standard LaTeX picture environment and requires no graphics package.
% ============================================================================

\begingroup
\setlength{\unitlength}{1mm}
\begin{picture}(145,62)
    % Main generation and verification path.
    \put(0,44){\framebox(23,12){\shortstack[c]{Prompt\\\(x\)}}}
    \put(23,50){\vector(1,0){8}}

    \put(31,44){\framebox(25,12){\shortstack[c]{Policy\\\(\pi_\theta\)}}}
    \put(56,50){\vector(1,0){8}}

    \put(64,44){\framebox(27,12){\shortstack[c]{Generated response\\\(y\sim\pi_\theta\)}}}
    \put(91,50){\vector(1,0){8}}

    \put(99,44){\framebox(22,12){\shortstack[c]{Verifier\\\(C(x,y)\)}}}
    \put(121,50){\vector(1,0){8}}

    \put(129,44){\framebox(16,12){\shortstack[c]{Reward\\\(v(x,y)\)}}}

    % Reference-policy and value-function paths.
    \put(5,18){\framebox(30,12){\shortstack[c]{Reference policy\\\(\pi_{\mathrm{ref}}\)}}}
    \put(35,24){\vector(1,0){7}}
    \put(42,18){\framebox(28,12){\shortstack[c]{KL penalty\\\(D_{\mathrm{KL}}\)}}}

    \put(82,18){\framebox(28,12){\shortstack[c]{Value model\\\(V_\psi(s_t)\)}}}
    \put(110,24){\vector(1,0){6}}
    \put(116,18){\framebox(29,12){\shortstack[c]{Advantage\\\(\widehat A_t\)}}}

    % Inputs to the PPO update.
    \put(56,18){\line(1,-1){6}}
    \put(62,12){\vector(1,0){18}}

    \put(130,18){\line(-1,-1){6}}
    \put(124,12){\vector(-1,0){4}}

    \put(137,44){\line(0,-1){16}}
    \put(137,28){\line(-1,0){37}}
    \put(100,28){\vector(0,-1){16}}

    \put(80,0){\framebox(40,12){\shortstack[c]{PPO update\\policy and value}}}
\end{picture}
\endgroup

    \caption{Conceptual structure of PPO-based reinforcement learning with
    verifiable rewards. The policy generates a response that is evaluated by
    a verifier. The resulting reward, the KL penalty relative to the reference
    policy, and advantage estimates derived from the value model inform the
    PPO update. Original schematic based on the RLVR framework described by
    \textcite{lambert2025tulu3}.}
    \label{fig:rlvr-training-pipeline}
\end{figure}

\subsection{Verifiable data, training recipe, and empirical behavior}
\label{subsec:rlvr-data-recipe-results}

An RLVR dataset must provide more than input prompts. Each prompt must be
associated with a verification procedure that can evaluate a generated
response. A convenient representation is

\begin{equation}
    \mathcal{D}_{\mathrm{RLVR}}
    =
    \left\{
        (x_i,C_i)
    \right\}_{i=1}^{N},
    \label{eq:rlvr-dataset}
\end{equation}

where \(x_i\) is a training prompt, \(C_i\) is its associated verifier, and
\(N\) is the number of prompts. The response itself is sampled from the
current policy during training; the dataset supplies the prompt and the
procedure for evaluating the response.

The verifier is task-dependent. In mathematical reasoning, it can extract
the final answer and compare it with a known answer. In precise instruction
following, it can test whether a completion satisfies a specified constraint.
The same binary reward interface can therefore be used across tasks whose
underlying verification procedures differ.

For GSM8K, Tülu 3 uses the training set of 7,473 prompts. Each example is
augmented with the standard eight-shot prompt used during evaluation to
encourage chain-of-thought generation. The verifier extracts the final
number from the generated response and compares it with the ground-truth
label. In abstract form,

\begin{equation}
    C_{\mathrm{GSM8K}}(x,y)
    =
    \mathbb{I}
    \left[
        E_{\mathrm{num}}(y)=a^\star(x)
    \right],
    \label{eq:rlvr-gsm8k-verifier}
\end{equation}

where \(E_{\mathrm{num}}\) is the answer-extraction procedure and
\(a^\star(x)\) is the ground-truth answer.

For MATH, the training mixture contains 7,500 prompts. The examples are
augmented with the standard three-shot chain-of-thought prompt used during
evaluation. The generated answer is extracted and checked using the
flexible MATH evaluation logic. This is important because mathematically
equivalent answers can have different textual representations. A verifier
based only on literal string equality could incorrectly reject two
representations of the same mathematical result. The MATH verification
procedure is therefore represented more generally as

\begin{equation}
    C_{\mathrm{MATH}}(x,y)
    =
    \mathbb{I}
    \left[
        E_{\mathrm{MATH}}(y)
        \equiv a^\star(x)
    \right],
    \label{eq:rlvr-math-verifier}
\end{equation}

where \(\equiv\) denotes the comparison logic used by the evaluation
procedure.

The instruction-following component is constructed differently. Tülu 3
randomly samples instructions from the Tülu 2 supervised fine-tuning mixture
and combines them with constraints drawn from the taxonomy used for IFEval.
A prompt-specific verifier is constructed for each constraint template.
If \(g_j(x,y)\) denotes the test associated with constraint template \(j\),
then the verification function can be represented as

\begin{equation}
    C_{\mathrm{IF}}(x,y)
    =
    \mathbb{I}[g_j(x,y)=1].
    \label{eq:rlvr-if-verifier}
\end{equation}

The condition \(g_j(x,y)=1\) indicates that the generated completion
satisfies the corresponding constraint. Unlike mathematical answer
verification, this procedure need not extract one numerical answer; it
checks the property specified by the instruction template.

The resulting Tülu 3 training mixture contains 29,946 prompts, divided
equally between mathematical tasks and verifiable instruction following.
The composition is shown in \cref{tab:rlvr-dataset-composition}.

\begin{table}[htbp]
    \centering
    \caption{Composition of the Tülu 3 verifiable-reward prompt dataset.
    Source: \textcite{lambert2025tulu3}.}
    \label{tab:rlvr-dataset-composition}
    \begin{tabularx}{\textwidth}{l r X}
        \toprule
        \textbf{Training domain}
        & \textbf{Prompts}
        & \textbf{Verification procedure}\\
        \midrule
        GSM8K
        & 7,473
        & Exact match against the extracted answer\\
        MATH
        & 7,500
        & Extracted-answer verification using flexible MATH evaluation logic\\
        Verifiable instruction following
        & 14,973
        & Prompt-specific constraint verifiers\\
        \midrule
        \textbf{Total}
        & \textbf{29,946}
        & \textbf{Binary verification}\\
        \bottomrule
    \end{tabularx}
\end{table}

The mathematical portion contains

\begin{equation}
    N_{\mathrm{math}}
    =
    7{,}473+7{,}500
    =
    14{,}973
    \label{eq:rlvr-math-count}
\end{equation}

prompts, which equals the 14,973 prompts in the verifiable instruction-
following component. The mixture therefore contains equal numbers of
mathematical and instruction-following prompts. Although the reward
interface is shared, the meaning of a successful response remains
domain-dependent because each task uses its own verifier.

The Tülu 3 procedure applies RLVR after supervised fine-tuning and
preference fine-tuning. The starting policy is therefore already an
instruction-tuned model that has undergone preference optimization. In
the reported experiments, the authors use DPO checkpoints as the principal
starting points and also investigate initialization from a weaker
supervised fine-tuning checkpoint. This ordering separates the roles of the
post-training stages. Supervised fine-tuning teaches the model from
demonstrations, preference fine-tuning optimizes relative response
preferences, and RLVR optimizes sampled responses according to verifiable
outcomes.

The PPO configuration used in the experiments is summarized in
\cref{tab:rlvr-ppo-hyperparameters}. The values distinguish the standard
configuration reported for RLVR from settings that vary with model scale or
the particular run.

\begin{table}[htbp]
    \centering
    \caption{Selected PPO hyperparameters reported for the Tülu 3 RLVR
    experiments. Values are transcribed from the RLVR configuration in
    \textcite{lambert2025tulu3}.}
    \label{tab:rlvr-ppo-hyperparameters}
    \small
    \begin{tabularx}{\textwidth}{l l X}
        \toprule
        \textbf{Hyperparameter}
        & \textbf{Reported value}
        & \textbf{Role or qualification}\\
        \midrule
        Discount factor \(\gamma\)
        & \(1.0\)
        & No geometric discounting of future rewards\\
        GAE parameter \(\lambda\)
        & \(0.95\)
        & Controls the temporal weighting in advantage estimation\\
        PPO clipping coefficient \(\epsilon\)
        & \(0.2\)
        & Clips the probability ratio to \([0.8,1.2]\) in the surrogate objective\\
        Value-function coefficient \(c_1\)
        & \(0.1\)
        & Weights the value-function loss\\
        Gradient-norm threshold
        & \(1.0\)
        & Limits the gradient norm\\
        Learning-rate schedule
        & Linear
        & Linear schedule during training\\
        Generation temperature
        & \(1.0\)
        & Sampling temperature for response generation\\
        Maximum prompt length
        & 2,048 tokens
        & Maximum prompt context length in the reported configuration\\
        Maximum token length
        & 2,048 tokens
        & Maximum token limit in the reported configuration\\
        Response length
        & 2,048 tokens
        & 1,024 tokens for GSM8K in the standard RLVR configuration\\
        Learning rate
        & \(3\times10^{-7}\)
        & \(1\times10^{-7}\) for the reported 70B configuration\\
        Effective batch size
        & 224
        & 640 for the reported 70B configuration\\
        PPO update iterations
        & 4
        & Number of PPO update iterations in the RLVR column\\
        Total episodes
        & 100,000
        & The final 70B run is reported separately with 400,000 episodes\\
        KL coefficient \(\beta\)
        & \(\{0.1,0.05,0.03,0.01\}\)
        & Sweep values in the standard RLVR configuration\\
        Warmup ratio \(\omega\)
        & \(0\) or \(0.1\)
        & Final 8B setting: \(0\); Table 21 70B setting: \(0.07\)\\
        Non-EOS penalty
        & \(-10\)
        & Additional penalty for responses without an EOS token\\
        \bottomrule
    \end{tabularx}
\end{table}

The paper reports a numerical inconsistency for the final 70B KL coefficient.
Table 21 gives \(\beta=0.07\), while the prose describing the final 70B
run in Section 6.4 reports \(\beta=0.7\). These values differ by a factor
of ten, and the paper does not reconcile them. They should therefore not be
silently treated as equivalent when reproducing the experiment. The same
final-run discussion reports a learning rate of \(1\times10^{-7}\), a
warmup ratio of \(0.1\), a response length of 2,048 tokens, 400,000
episodes, and an effective batch size of 640. The warmup value also differs
from the \(0.07\) listed in Table 21, so the prose and table should be
consulted separately when reconstructing the exact run
\parencite{lambert2025tulu3}.

Several implementation decisions are important for interpreting these
settings. The value model is initialized from a general reward model rather
than directly from the DPO policy. The general reward model is trained on
UltraFeedback data, but its predicted score is not used as the primary task
reward in RLVR. It initializes the value model, whose purpose is to estimate
returns for PPO. This distinction is important: replacing a learned reward
model with a verifier does not remove the need for a value function in an
actor--critic implementation.

The authors disable dropout during reward-model and RL training. As
explained above, PPO compares token log probabilities obtained during
rollout and learning. If dropout changes the forward computation between
these phases, the probability ratio may not equal one at the start of the
first update. This can lead to excessive clipping and weak or vanishing
policy gradients. Disabling dropout makes the probability calculations
consistent under the same model parameters and inputs.

The available prompt set is reused across multiple training epochs. In one
GSM8K ablation, approximately 100,000 episodes are run using 7,473 prompts,
which corresponds to

\begin{equation}
    \frac{100{,}000}{7{,}473}
    \approx 13.38
    \label{eq:rlvr-prompt-reuse}
\end{equation}

passes over the prompt set if reuse is approximately uniform. The prompts
are shuffled between epochs. The authors also evaluate checkpoints during
training and select checkpoints using development-set performance rather
than relying only on the final training reward.

A simple numerical example clarifies the expected-reward interpretation.
Suppose, hypothetically, that a policy samples eight responses to one prompt
and that five pass the verifier. With \(\alpha=10\), the empirical success
frequency is \(5/8=0.625\), and the empirical mean verification reward is

\begin{equation}
    \widehat{\mathbb{E}}[v]
    =
    \frac{5(10)+3(0)}{8}
    =
    6.25.
    \label{eq:rlvr-worked-reward}
\end{equation}

This value is an estimate of the expected reward for that prompt under the
sampling policy. It is not a direct estimate of general reasoning ability:
it measures the probability of satisfying the particular verifier under the
particular prompt and generation procedure. A different verifier or prompt
distribution can yield a different expected reward.

The Tülu 3 experiments examine several factors that influence the final
model. Initializing the value model from a general reward model performs
better in the reported GSM8K ablation than initializing it from the DPO
policy. Starting RLVR from a weaker SFT model can reach a similar level of
training verification reward, but it incurs greater KL divergence than
starting from the stronger DPO model; the authors also find that a stronger
starting model generally produces better test performance. These findings
show that reaching a particular training reward and generalizing well are
not identical outcomes.

The authors also test whether reward-model scores should be added to
verifiable rewards. In their GSM8K experiments, using the verifiable reward
alone performs better than combining it with reward-model scores, while the
combined signal is noisier. This result supports the use of a direct
verification signal when a reliable verifier is available, rather than
assuming that an additional learned score must improve the optimization
signal.

The final performance comparison is reported for the 8B and 70B Tülu 3
models. Selected results are reproduced in
\cref{tab:rlvr-final-performance}. The comparison is between the DPO
starting point and the RLVR-trained model; the base instruction-tuned model
is included in the source paper but is not needed to isolate the effect of
the RLVR stage.

\begin{table}[htbp]
    \centering
    \caption{Selected benchmark results for the Tülu 3 DPO starting points
    and RLVR-trained models. Scores are reproduced from Table 23 of
    \textcite{lambert2025tulu3}.}
    \label{tab:rlvr-final-performance}
    \small
    \begin{tabularx}{\textwidth}{l r r r r}
        \toprule
        \textbf{Evaluation}
        & \textbf{8B DPO}
        & \textbf{8B RLVR}
        & \textbf{70B DPO}
        & \textbf{70B RLVR}\\
        \midrule
        Average score
        & 64.4 & 64.8 & 75.9 & 76.0\\
        MATH
        & 42.0 & 43.7 & 62.3 & 63.0\\
        GSM8K
        & 84.3 & 87.6 & 93.5 & 93.5\\
        IFEval (strict)
        & 81.1 & 82.4 & 82.6 & 83.2\\
        BigBenchHard
        & 65.8 & 66.0 & 81.8 & 82.0\\
        DROP
        & 62.5 & 62.6 & 74.1 & 74.3\\
        Safety6 task average
        & 87.2 & 85.5 & 89.0 & 88.3\\
        \bottomrule
    \end{tabularx}
\end{table}

At the 8B scale, the average score increases from 64.4 to 64.8. The
GSM8K score rises by 3.3 percentage points, from 84.3 to 87.6, while MATH
increases by 1.7 points and strict IFEval by 1.3 points. These changes
demonstrate targeted gains on the mathematical and instruction-following
tasks used to construct the verifiable training mixture. The average
improvement is nevertheless modest, and not every evaluation improves:
the Safety6 task average decreases from 87.2 to 85.5.

At the 70B scale, the average score increases from 75.9 to 76.0. MATH
increases from 62.3 to 63.0, and strict IFEval increases from 82.6 to
83.2, while GSM8K remains at 93.5. The authors suggest that the lack of
GSM8K improvement is related to the strong initial score, which leaves
less room for improvement on that benchmark. The Safety6 task average
also decreases slightly, from 89.0 to 88.3. The results therefore do not
support the claim that RLVR improves every capability of a general-purpose
language model. They support the narrower conclusion that it can improve
selected capabilities while other dimensions remain unchanged or decline.

The implementation also has substantial computational requirements.
The Tülu 3 infrastructure uses ZeRO Stage 3 to distribute model parameters
and applicable optimizer states across devices. The RLVR setup contains a
policy model, a reference policy model, and a value model. The policy and
value models are trained, whereas the reference policy is used for
inference. The authors allocate dedicated inference GPUs and use Ray for
resource allocation and vLLM with PagedAttention for efficient generation.
PagedAttention uses paging-based memory management to reduce memory
fragmentation and redundant duplication, enabling larger inference batches.

The authors further use asynchronous reinforcement learning. In a
synchronous setup, inference GPUs first generate rollout data and learner
GPUs subsequently optimize the policy. These stages can leave one set of
GPUs idle while the other is working. The Tülu 3 implementation overlaps
inference and training to reduce this idle time. Asynchronous generation,
however, can produce stale trajectories if inference generates data faster
than the learner consumes it. To improve reproducibility, the implementation
trains on the second most recent inference data rather than arbitrarily old
rollouts. The reported RL run times are approximately 65 hours for the 8B
run on 8 GPUs and 60 hours for the 70B run on 48 GPUs
\parencite{lambert2025tulu3}. These figures are implementation-specific
measurements, not universal estimates of the cost of RLVR.

\subsection{Limitations, interpretation, and connection to GRPO}
\label{subsec:rlvr-limitations-transition}

The first limitation of RLVR is its dependence on the quality of the
verification function. The policy is optimized against the criterion
implemented by the verifier, not against an abstract notion of correctness
that is independent of that implementation. If the verifier captures the
intended objective reliably, maximizing its expected reward can improve the
corresponding capability. If the verifier is incomplete or incorrectly
specified, the model can learn to satisfy the implemented condition without
achieving the intended outcome.

For example, a verifier that checks only whether a required phrase appears
in a response can reward a completion containing that phrase even if the
response does not provide the explanation that the task requires. In such
a case, policy optimization can be successful with respect to the formal
reward and unsuccessful with respect to the intended task. This is not a
failure specific to PPO; it follows from the definition of the optimization
objective. The verifier defines the outcome being optimized, while the
optimizer determines how the policy is changed to increase that outcome.

A second limitation follows from the sparse binary reward. In a mathematical
task, a response with largely correct intermediate reasoning but an
incorrect final answer can receive the same zero reward as a response with
entirely incorrect reasoning. Conversely, two correct responses receive
the same positive reward even if one is more concise or uses a more
efficient derivation. The mapping

\begin{equation}
    v(x,y)\in\{0,\alpha\}
    \label{eq:rlvr-sparse-reward}
\end{equation}

discards information about degrees of partial correctness, reasoning quality,
and the efficiency of the solution. The value function and advantage
estimator can improve the statistical use of sampled trajectories, but
they do not add information that is absent from the verifier's reward.

A third limitation is overoptimization. Reducing \(\beta\) makes divergence
from the reference policy less costly in the objective and generally permits
the policy to move farther from its initial distribution. This can increase
the targeted verification reward, but it can also damage performance on
other evaluations or produce undesirable forms of specialization. The
Tülu 3 experiments report that higher KL divergence is often associated
with lower average scores in some settings, although this relationship is
not universal. The authors also identify overoptimized outputs in some
higher-KL instruction-following runs. Thus, KL divergence is a useful
diagnostic, but neither a large nor a small value should be interpreted
without considering task performance and independent evaluations.

A fourth limitation concerns generalization. The training objective is
defined over a finite distribution of prompts and associated verifiers.
High reward on this distribution does not mathematically guarantee transfer
to new prompts, new task formulations, or capabilities not measured by the
verifier. The Tülu 3 results show that training improvements can translate
to corresponding test benchmarks, but they also show that broader evaluation
metrics may remain unchanged or decline. Training reward must therefore be
considered together with held-out evaluations and broader capability
measurements.

The experiments also show that a verifier-based reward is not automatically
improved by combining it with a learned reward-model score. In the reported
GSM8K ablation, the verifiable reward alone performs better, while the
combined signal is noisier. This is a useful methodological result, but it
should not be generalized to every possible task or verifier. The conclusion
supported by the experiment is that adding a learned score is not inherently
beneficial when a direct verification signal is already available.

Finally, PPO-based RLVR has nontrivial computational demands. Response
generation, policy optimization, reference-policy evaluation, and value
estimation must be coordinated. Distributed model placement and asynchronous
rollouts can improve hardware utilization, but asynchronous execution
introduces the possibility of stale data and reproducibility difficulties.
The training recipe is consequently part of the method: a mathematically
appropriate reward objective is not sufficient if its implementation
produces unstable probability ratios, unreliable advantage estimates, or
rollouts that are too stale for the intended update.

The central result of this section is that verifiable outcomes can provide
a direct reinforcement-learning signal for language models whenever a
reliable evaluator exists. RLVR retains the autoregressive policy, uses a
binary or otherwise explicitly verifiable reward, regularizes policy
movement relative to a reference distribution, and employs an optimizer
such as PPO to convert sampled outcomes into parameter updates. The Tülu 3
experiments demonstrate the approach on mathematical reasoning and precise
instruction following, while also exposing the importance of verifier
quality, value estimation, KL control, initialization, generalization, and
computational infrastructure.

This formulation separates the definition of the reward from the choice of
the policy-optimization algorithm. PPO uses a learned value function to
construct advantages and a clipped probability-ratio objective to control
updates. The next section examines Group Relative Policy Optimization
(GRPO), an alternative optimization approach for verifiable-reward
training. It builds on the same general setting of sampled responses and
verifiable outcomes but changes how relative performance information is
used to construct the policy update.
